\Needspace{24\baselineskip}
\chapterrule
\section{系统已经会用工具了，怎样稳定部署并控制成本？}
\label{intro:23b-deployment-and-ecosystem}
\chaptercompass{模型能使用工具之后，部署还受什么制约？}{分清权重存储、KV 缓存、量化与实际服务的成本。}{$N$ 个参数 $\to$ 权重字节；$T$ 个 Token $\to$ KV 缓存}{本节推理使用固定权重；缓存和运行状态会变化}
\modelpipelinefocus{8}{部署、缓存与系统}


Agent 的报告流程已经包含检索、计算与人工复核，
但模型仍要在现实硬件上运行，且必须考虑缓存、延迟、成本与隐私。
这一章转向基础设施，不把性能优化误写成模型在对话中重新学习。

\subsection{训练代码、模型权重、推理程序是三件事}

让一个大语言模型工作，至少需要区分：
\begin{enumerate}
\item \textbf{模型结构与代码}：决定矩阵和非线性如何排列、计算。
\item \textbf{模型权重}：训练得到的大量参数数值。
\item \textbf{推理运行时}：加载权重、完成 Tokenization、执行网络计算并生成输出的程序。
\end{enumerate}
有结构和代码但没有训练好的权重，模型还无法提供成熟的语言能力；
拿到权重后，也需要合适的推理程序才能运行。

\textbf{还要区分“公司／模型／产品”三个层次。}
\textbf{公司}是研发或运营的组织；
\textbf{模型}是有明确架构、参数和版本的计算系统；
\textbf{产品}是用户实际接触的应用、界面或 API，
往往还组合检索、工具调用、审核与计费。
同一模型可服务多个产品；同一产品也可接入多个模型。
因此比较产品与模型时，要分别记录产品功能、使用的权重版本和模型配置。

\textbf{开放权重}意味着外部能够获取模型参数；开源程度则应分别核对训练代码、数据、再分发条款与商业许可。

\textbf{私有化部署}一般是指在自有设备或受控制的基础设施上运行模型。它有助于控制数据流向，但并不自动带来安全保障：仍需考虑访问权限、日志、模型许可和升级流程。

\subsection{模型为什么需要显存？}

权重由许多数字组成。假设模型有 $10^9$（十亿）个参数，如果每个参数用 2 字节存储，仅权重就占约
\[
10^9\times2\ \text{bytes}\approx2\ \text{GB}.
\]
这是十进制粗略估计，还不包含 \conceptlink{kv-cache}{KV Cache}、中间激活、运行时缓冲、框架开销等。

\textbf{GPU} 擅长对大量数字并行计算；\textbf{TPU} 是面向机器学习工作负载设计的专用加速器；\textbf{NPU} 泛指神经网络计算加速单元。这些硬件为模型提供执行矩阵运算和数据搬运的计算基础设施。

\begin{warningbox}
训练时还常需保存 \conceptlink{gradient}{梯度}、优化器状态等额外数据，这些都要在权重存储之外另计显存。即使模型权重能放入显存，运行长上下文或高并发时仍可能发生内存不足。

硬件名称说明设备类别；算子与精度支持、带宽和延迟须按具体型号、运行时与负载核对。
\end{warningbox}


\subsection{为什么大模型训练要拆到多张卡？}

权重只是训练显存的一部分。还要保存梯度、优化器状态和
反向传播需要的激活。以十亿参数为教学估算：
仅 FP32 主权重、梯度与 Adam 的两份状态，
理论上就需要大约 $4\times4=16$ GB，
激活、临时缓冲和额外权重副本需另计。
显存不足时需要\textbf{分布式训练}与分片。

\begin{center}\small
\tagpdfsetup{table/header-rows={1}}
\begin{tabularx}{\linewidth}{@{}>{\raggedright\arraybackslash}p{35mm}>{\raggedright\arraybackslash}X@{}}
\toprule
做法 & 分摊的对象及主要代价\\\midrule
数据并行（Data Parallelism） & 不同设备处理不同样本批次，再同步或归约梯度；朴素方案每卡仍保留完整权重\\
张量并行（Tensor Parallelism） & 把一层的大矩阵与计算拆到多设备；层内通信频繁\\
流水线并行（Pipeline Parallelism） & 不同设备保存不同层；微批次填充流水线，需处理流水线空泡\\
状态分片（如 ZeRO/FSDP） & 分摊优化器状态、梯度乃至权重；以设备间通信换显存\\
\bottomrule
\end{tabularx}
\end{center}

\textbf{混合精度训练（Mixed Precision）}用如 BF16 的较短浮点表示
执行适合的矩阵计算，并对某些归约、统计或状态保留较高精度。
BF16 具有与 FP32 相同宽度的指数位，因此动态范围较大，
但尾数精度更低；采用哪种精度取决于硬件及数值稳定性。
\textbf{FlashAttention} 则通过分块、在线 Softmax 和减少高带宽显存读写
提高标准 Attention 的计算效率；它在适用条件下
\textbf{保持相同的注意力数学目标}，不是一种新的注意力训练损失。
工程选型要同时测量吞吐、显存、通信、数值误差与最终模型质量。

\subsection{量化：为什么更少的位数可以降低成本？}
\phantomsection\label{concept:quant}

原来每个权重使用较高数值精度保存；\textbf{量化}用更少位数表示近似的权重或激活。例如从每个参数 16 bit 变为 8 bit，权重存储理论上可下降到原来的一半左右。

\begin{example}
\noindent\emph{把四个权重实际量化到 8 bit}\par

取浮点权重 $(-1,-0.2,0.2,1)$，
使用包含零点的\textbf{对称有符号 8-bit 量化}，
整数范围 $[-127,127]$，比例尺 $s=1/127$。
将每个数除以 $s$、四舍五入并截取合法范围：
\[
(-1,-0.2,0.2,1)
\longrightarrow(-127,-25,25,127).
\]
反量化（整数乘回 $s$）得到约
$(-1,-0.197,0.197,1)$。
存储从每数 16 bit 变为每数 8 bit 时，
仅这四个权重从 8 字节降为 4 字节，不包括比例尺和分组元数据。
\end{example}

代价是近似误差：相近权重会落入同一个量化格。模型效果是否受影响，取决于量化方案、模型分布和任务，并不是位数越少越好。

\subsection{蒸馏、剪枝、LoRA 不是同一种事}
\phantomsection\label{concept:lora}

\begin{center}
\small
\renewcommand{\arraystretch}{1.23}
\tagpdfsetup{table/header-rows={1}}
\begin{tabularx}{\linewidth}{@{}p{38mm}X@{}}
\toprule
技术 & 核心目的\\
\midrule
量化（Quantization） & 用低精度数值减少模型存储或计算成本\\
蒸馏（Distillation） & 利用较大或较强模型的输出作为训练信号，帮助学生模型学习\\
剪枝（Pruning） & 去除或稀疏化较不重要的权重、连接或结构\\
LoRA & 在微调阶段用低秩增量参数，减少需要训练的参数量\\
\bottomrule
\end{tabularx}
\end{center}

\begin{example}
\noindent\emph{硬标签与软标签怎样不同？}\par

教师模型给两个类别的概率为 $(0.8,0.2)$。
硬标签只告诉学生“类别 1 是正确答案”，可表示成 $(1,0)$；
\textbf{软标签}则告诉学生：类别 2 仍有 $0.2$ 的概率。
如果学生预测 $(0.6,0.4)$，硬标签交叉熵为
$-\log0.6\approx0.511$，
而对教师软标签的交叉熵为
\[
-\bigl(0.8\log0.6+0.2\log0.4\bigr)\approx0.592.
\]
后者同时考虑两个类别的匹配程度。蒸馏时还可以调整温度，并把软标签损失与原任务损失结合。
\end{example}

\begin{example}
\noindent\emph{LoRA 的参数量真正省在哪里？}\par

对一个 $1000\times1000$ 的方阵，完整训练它需要
$1000^2=1{,}000{,}000$ 个可学习权重。
LoRA 若冻结原矩阵，用秩 $r=8$ 的增量
$\Delta W=AB$，$A:1000\times8$，
$B:8\times1000$，则需要
\[
1000\times8+8\times1000
=\boxed{16{,}000}
\]
个可学习参数，约为完整权重的 $1.6\%$。
这指的是该矩阵新训练的增量参数数目，
不是整套模型参数或部署占用必然缩小到 $1.6\%$。
\end{example}

这些技术解决的成本不同：量化和剪枝常用于减小部署开销，LoRA 主要减少微调时需要更新的参数。

\subsection{从文本模型到多模态模型}

\textbf{多模态}是模型能处理文字、图片、音频、视频等多种类型输入或输出。具体模型可使用模态专用编码器或共享表示；进入语言主干后，再用注意力、投影等机制融合信息。

常见缩写：
\begin{itemize}
\item \textbf{ASR}（Automatic Speech Recognition）：语音识别，把声音转成文字。
\item \textbf{TTS}（Text to Speech）：将文字转换为合成语音。
\item \textbf{AIGC}（AI-Generated Content）：由 AI 生成或辅助生成内容的广义说法。
\item \textbf{AGI}（Artificial General Intelligence）：通用人工智能，是仍在讨论的目标与概念，具体主张需要给出测试定义。
\end{itemize}

\subsection{为什么还有框架、镜像、推理引擎和协议？}

模型权重之外仍需要很多工程工具。一个典型流程是：
\[
\boxed{\text{模型权重}\longrightarrow
\text{推理引擎}\longrightarrow
\text{接口服务}\longrightarrow
\text{应用与工具}}
\]
\textbf{框架}用于定义和执行训练/推理计算；\textbf{镜像}打包运行环境；\textbf{推理引擎}专注批处理、缓存和计算性能；\textbf{接口}让其他程序以约定的格式请求模型服务；\textbf{协议}帮助多种工具和系统建立协作规则。

\textbf{API、\conceptlink{agent}{Agent} 和 \conceptlink{mcp}{MCP} 等概念主要属于系统集成层。}工具调用通过上下文和外部动作扩展产品能力，本书的推理流程使用固定模型参数。

\subsection{技术生态地图：CUDA、PyTorch、Hugging Face、Ollama、vLLM}

同一个模型可以通过不同软件组合训练、下载和部署。
不要把\textbf{硬件计算平台、训练框架、模型仓库、运行工具和推理服务}
当作可以相互替代的五种模型架构。

\begin{center}
\small
\renewcommand{\arraystretch}{1.2}
\tagpdfsetup{table/header-rows={1}}
\begin{tabularx}{\linewidth}{@{}>{\raggedright\arraybackslash}p{26mm}>{\raggedright\arraybackslash}p{34mm}>{\raggedright\arraybackslash}X@{}}
\toprule
项目 & 主要定位 & 在一条常见工作链中做什么\\\midrule
CUDA & NVIDIA GPU 计算平台与工具链 &
提供 GPU 并行计算接口、编译器与加速库；本身不是神经网络\\
PyTorch & 张量计算与自动求导框架 &
定义模型、训练/反传，也可执行推理；可使用 CPU 或支持的加速设备\\
Hugging Face & 模型、数据与工具生态 &
Hub 管理模型/数据仓库和说明文档；Transformers 等库加载和运行模型\\
Ollama & 本地模型管理与运行工具 &
在支持的机器上拉取、运行模型，并提供本地服务接口\\
vLLM & 面向服务的模型推理引擎 &
组织高并发批处理、KV Cache 和模型服务；强调推理吞吐与资源利用\\
\bottomrule
\end{tabularx}
\end{center}

例如，开发者从 Hugging Face Hub 取得一份\textbf{允许使用}的权重，
用 PyTorch 完成研究实验，再选合适的推理运行时对外提供服务。
若选用 NVIDIA GPU，底层可使用 CUDA；也可以按照设备和模型支持情况
选择其他计算后端。Ollama 偏向便捷的本地运行，
vLLM 偏向服务型推理；是否更快，要在相同的权重、硬件、
输入长度、并发数和量化条件下实际测试。

\textbf{查询入口：}
\href{https://docs.nvidia.com/cuda/cuda-programming-guide/index.html}{CUDA}、
\href{https://docs.pytorch.org/docs/stable/index.html}{PyTorch}、
\href{https://huggingface.co/docs/hub/index}{Hugging Face Hub}、
\href{https://docs.ollama.com/}{Ollama}、
\href{https://docs.vllm.ai/}{vLLM}。
这些工具的\textbf{接口、兼容硬件、部署模式和许可证可能更新}。
本表只定位各层责任，不固定特定版本、吞吐数字或产品优劣排名。

\chapterchecktasks{23b-deployment-and-ecosystem}{正文基础三题}{
\item \textbf{区分：}训练代码、模型权重、推理引擎和对外 API 在一个产品里承担什么不同职责？\chapterinlinehint{区分训练程序、权重、推理引擎与访问 API。}
\item \textbf{估算：}$10^9$ 个参数各占 $2$ 字节，不计运行开销仅权重需要多少十进制 GB？若另有 $0.5$ GB 缓存，合计至少多少？\chapterinlinehint{权重 $2$ GB；加给定缓存至少 $2.5$ GB。}
\item \textbf{找错：}有人认为把参数量化到更少位数，模型质量、吞吐与延迟都会按字节比例等幅改善。请指出两处需要实测的地方。\chapterinlinehint{字节减少不直接保证质量或端到端加速。}
}
\chapterrecallanswer{权重、KV Cache 与运行时开销是不同内存项；量化节省存储不等于任务质量和时延自动同比变化。}



\begin{conclusion}
部署一个 AI 产品需要模型结构、训练权重、推理引擎、接口与安全措施共同工作。
模型的计算能力与系统的运行能力相互配合，才能构成可用的服务。
接下来再回头审视经典 Transformer 自身的表示、计算与缓存瓶颈，
分析现代架构为什么会变化。
\end{conclusion}
